\documentclass[11pt]{article}
\usepackage[T1]{fontenc}
\usepackage[utf8]{inputenc}
\usepackage{amsmath,amssymb,mathtools}
\usepackage{booktabs,geometry,microtype}
\usepackage{xcolor,hyperref}
\geometry{margin=25mm}
\hypersetup{colorlinks=true,linkcolor=blue!50!black,urlcolor=blue!50!black}
\setlength{\emergencystretch}{2em}
\newcommand{\E}{\mathbb{E}}
\newcommand{\Cov}{\operatorname{Cov}}
\newcommand{\diag}{\operatorname{diag}}
\newcommand{\tr}{\mathsf T}
\newcommand{\Ic}{I_{\mathrm c}}
\newcommand{\Imis}{I_{\mathrm m}}
\newcommand{\Io}{I_{\mathrm o}}

\title{Natural Coordinates for GPURec:\\Diagnosis and Practical Optimization Suggestions}
\author{A standalone mathematical and experimental assessment}
\date{6 September 2026}

\begin{document}
\maketitle

\begin{abstract}
The gene-wise duplication--loss--transfer model admits a useful hierarchical
logit parameterization: one coordinate describes duplication/transfer intensity,
one their relative mixture, and one loss relative to speciation. For fixed
complete-history counts, including the appropriate survival-conditioning
augmentation, its objective is a sum of three scalar binary likelihoods with an
exactly diagonal Hessian. This does not diagonalize the observed likelihood:
uncertainty over reconciliation histories reintroduces coupling and weak
curvature. The most defensible practical strategy is therefore to use
complete-history counts for a short bounded EM initialization and a cheap
curvature seed, then learn the observed curvature with safeguarded per-family
BFGS updates. Original log-rate bounds and convergence tests should remain
authoritative. This note separates exact identities, implementation suggestions,
and the experimental evidence available at the time of writing.
\end{abstract}

\section{Diagnosis: the coordinates are only part of the difficulty}

For one family, write
\[
 \theta=(\theta_D,\theta_L,\theta_T)
       =(\log_2 D,\log_2 L,\log_2 T),\qquad
 Z=1+D+L+T.
\]
Speciation is the reference event, with rate one, so
\begin{equation}
 (p_S,p_D,p_L,p_T)=\frac{(1,D,L,T)}{Z}.
 \label{eq:probabilities}
\end{equation}
This discussion holds recipient and other non-rate parameters fixed, as in the
gene-wise Coleman fits. Their contributions to each complete history are then
independent of these three rates.

The existing log rates are already canonical exponential-family parameters, up
to multiplication by $\ln 2$. Thus they are not an intrinsically inappropriate
starting point. For fixed nonnegative expected event counts $N_k$, put
$N=\sum_k N_k$. The complete-history negative log-likelihood surrogate in bits is
\begin{equation}
 Q_\theta(\theta;N)
 =N\log_2 Z-N_D\theta_D-N_L\theta_L-N_T\theta_T+\text{constant}.
 \label{eq:surrogate-theta}
\end{equation}
It is convex, with
\begin{equation}
 \nabla Q_\theta=Np-(N_D,N_L,N_T)^\tr,\qquad
 C_\theta(\theta;N)=\ln(2)N\{\diag(p)-pp^\tr\},
 \quad p=(p_D,p_L,p_T)^\tr.
 \label{eq:surrogate-curvature}
\end{equation}

The observed family does not specify its event counts. Alternative histories can
explain similar gene-tree patterns using different combinations of duplication,
transfer, loss, and extinct lineages. Moving the rates changes that posterior
mixture. Consequently, observed curvature can be much weaker than complete-data
curvature, strongly coupled, and indefinite away from a solution. This is a
statistical ambiguity as well as a numerical scaling problem.

\section{An exact hierarchical coordinate system}

\subsection{Definition, interpretation, and inverse}

Define $\psi=(u,v,w)$ by
\begin{equation}
 \boxed{\quad
 u=\log_2\frac{D+T}{1+L},\qquad
 v=\log_2\frac{T}{D},\qquad
 w=\log_2 L.
 \quad}
 \label{eq:hierarchy}
\end{equation}
Here $u$ compares the event groups $\{D,T\}$ and $\{S,L\}$; $v$ compares
transfer with duplication within the first group; and $w$ compares loss with
speciation within the second. This directly exposes the observed
duplication--transfer substitution rather than mixing it into two absolute rates.

Let $F(x)=\log_2(1+2^x)$. The inverse is smooth for positive rates:
\begin{equation}
 \theta_D=u+F(w)-F(v),\qquad
 \theta_L=w,\qquad
 \theta_T=u+F(w)+v-F(v).
 \label{eq:inverse}
\end{equation}
Equivalently, $D=2^u(1+2^w)/(1+2^v)$, $L=2^w$, and $T=2^vD$.
Evaluate $F$ using a stable softplus implementation.

For $\sigma(x)=1/(1+e^{-x})$, set
\[
 a=\sigma(u\ln 2),\qquad b=\sigma(v\ln 2),\qquad
 c=\sigma(w\ln 2).
\]
Then the four-way event choice factors into binary choices:
\begin{equation}
 p_D=a(1-b),\qquad p_T=ab,\qquad
 p_S=(1-a)(1-c),\qquad p_L=(1-a)c.
 \label{eq:binary-factorization}
\end{equation}
This is a computationally natural factorization, not a claim that the original
log rates cease to be the canonical parameters of the model.

\subsection{Exact separation of the fixed-count surrogate}

Write $N_A=N_D+N_T$ and $N_B=N_S+N_L$, so $N=N_A+N_B$. Substituting
\eqref{eq:binary-factorization} into the surrogate gives
\begin{equation}
 Q_\psi(\psi;N)
 =\underbrace{NF(u)-N_Au}_{\text{group choice}}
 +\underbrace{N_AF(v)-N_Tv}_{\text{transfer versus duplication}}
 +\underbrace{N_BF(w)-N_Lw}_{\text{loss versus speciation}}
 +\text{constant}.
 \label{eq:surrogate-hierarchy}
\end{equation}
All counts in this differentiation are held fixed. In particular,
\begin{align}
 \nabla_\psi Q_\psi
 &=\begin{pmatrix}Na-N_A\\N_Ab-N_T\\N_Bc-N_L\end{pmatrix},
 \label{eq:hierarchical-gradient}\\
 C_\psi(\psi;N)
 &=\ln(2)\diag\!\left(Na(1-a),\;N_Ab(1-b),\;N_Bc(1-c)\right).
 \label{eq:hierarchical-curvature}
\end{align}
The Hessian is exactly diagonal. When all relevant counts are positive, the
unconstrained minimizer is
\begin{equation}
 u^+=\log_2\frac{N_A}{N_B},\qquad
 v^+=\log_2\frac{N_T}{N_D},\qquad
 w^+=\log_2\frac{N_L}{N_S}.
 \label{eq:hierarchical-em}
\end{equation}
After inversion, these are precisely $D^+=N_D/N_S$, $L^+=N_L/N_S$, and
$T^+=N_T/N_S$. Thus an exact unconstrained EM update is the same operation in
either coordinate system; a mere invertible reparameterization does not
accelerate that exact update. Zero counts and the actual rate bounds require
the constrained calculation below.

\section{Survival conditioning and missing information}

\subsection{Use genuine augmented counts, not an arbitrary signed score decomposition}

Let $A_y$ be the unconditioned likelihood numerator, $q$ the survival probability,
and $e=1-q$ the extinction probability. The fitted likelihood is $A_y/q$. Since
\[
 \frac{A_y}{q}=A_y\sum_{m=0}^{\infty}e^m,
\]
it admits an augmentation by the observed family and a geometrically distributed
number of extinct ``ghost'' families. For this augmentation,
\begin{equation}
 N_k=\E[n_k\mid y,\theta]
       +\frac{e}{q}\E[n_k\mid\text{extinct},\theta]\geq0.
 \label{eq:ghost-counts}
\end{equation}
These are the appropriate expected sufficient statistics for the EM surrogate.
Data-minus-survival expressions can yield the same observed score, but score
agreement alone does not establish an EM lower-bound argument for those signed
quantities.

The existing adjoint calculation can expose the four independent event
log-probability derivatives before the softmax chain rule:
\begin{equation}
 a_k=\frac{\partial\ell}{\partial\log_2 p_k}=-N_k,
 \qquad
 g_k=a_k-p_k\sum_j a_j=Np_k-N_k,
 \quad k\in\{D,L,T\}.
 \label{eq:count-check}
\end{equation}
The extinction and survival-normalization derivatives must be included. Transfer
recipient contributions are summed to obtain the family-level transfer count.
This makes count extraction approximately a by-product of the ordinary reverse
pass, not a second reconciliation solve.

\subsection{What diagonal complete-data curvature does not imply}

For an augmented complete history $h$, let
$s_\psi(h)=\nabla_\psi\log_2 P(h,y\mid\psi)$. With posterior expectations
evaluated at the current point, the observed Hessian in bits satisfies
\begin{equation}
 H_\psi=I_{\mathrm c,\psi}-I_{\mathrm m,\psi},\qquad
 I_{\mathrm c,\psi}=C_\psi(\psi;\E[n\mid y,\psi]),\qquad
 I_{\mathrm m,\psi}=\ln(2)\Cov[s_\psi(h)\mid y,\psi].
 \label{eq:missing-information}
\end{equation}
The covariance generally has off-diagonal entries. Diagonalizing the first term
therefore neither diagonalizes the observed Hessian nor removes weakly
identified directions. A diagonal complete-data metric can also be much too
stiff: the subtraction in \eqref{eq:missing-information} can be substantial.

For clarity, a nonlinear Hessian transformation is not just a change of basis.
If $J=\partial\theta/\partial\psi$ and $g=\nabla_\theta\ell$, then
\begin{equation}
 H_\psi=J^\tr H_\theta J
       +\sum_k g_k\nabla_\psi^2\theta_k.
 \label{eq:hessian-transform}
\end{equation}
For \eqref{eq:inverse}, the Jacobian and correction are
\begin{align}
 J&=\begin{pmatrix}1&-b&c\\0&0&1\\1&1-b&c\end{pmatrix},\\
 \sum_k g_k\nabla_\psi^2\theta_k
 &=\ln(2)(g_D+g_T)\diag\!\left(0,-b(1-b),c(1-c)\right).
 \label{eq:explicit-correction}
\end{align}
This gradient-dependent term explains why the hierarchy can remove much apparent
negative curvature away from an optimum. At an interior stationary point it
vanishes, and Hessian inertia is preserved under the remaining congruence.

At a regular interior EM fixed point, the local EM convergence matrix is
$R=\Ic^{-1}\Imis$. Its eigenvalues, equivalently those of
$\Ic^{-1/2}\Imis\Ic^{-1/2}$, are unchanged by a smooth invertible
reparameterization of the same exact EM map. A coordinate change can improve
numerical scaling, but cannot remove the underlying missing-information
fraction. These stationary-point statements should not be applied to an
arbitrary off-stationary Hessian. At a constrained optimum the relevant
second-order analysis uses feasible directions and the constraint geometry.

\section{Bounds are a substantive part of the problem}

Coleman uses $r_{\min}=10^{-6}$ and $r_{\max}=2$. Let
$\ell_b=\log_2 r_{\min}$ and $h_b=\log_2 r_{\max}$. The original box becomes
\begin{align}
 \ell_b&\leq w\leq h_b,\\
 \ell_b&\leq u+F(w)-F(v)\leq h_b,\\
 \ell_b&\leq u+F(w)+v-F(v)\leq h_b.
 \label{eq:coupled-bounds}
\end{align}
Replacing this region by three independent coordinate intervals changes the
optimization problem. Mapping a proposal back and clipping its rates preserves
feasibility, but changes the actual step; predicted decrease and acceptance
must be evaluated consistently with that step.

The bounded EM M-step is particularly cheap in the original log rates. For each
assignment of a coordinate to its lower bound, free state, or upper bound,
let $B$ be the pinned coordinates and $b_j$ their log-rate bounds. Set
\[
 c_B=1+\sum_{j\in B}2^{b_j},\qquad
 A_B=N_S+\sum_{j\in B}N_j.
\]
The free stationary solution is
\begin{equation}
 \theta_k^+=\log_2\frac{N_kc_B}{A_B},\qquad k\notin B.
 \label{eq:bounded-em}
\end{equation}
Enumerate the $3^3=27$ active-set assignments and check feasibility and KKT signs.
For the maximizing score $s_k=N_k-Np_k$, the signs are $s_k\leq0$ at a
lower bound, $s_k=0$ when free, and $s_k\geq0$ at an upper bound. This avoids
irreversible-pinning mistakes: pinning one rate can require releasing another.
Positive-total counts and finite bounds make the surrogate strictly convex
in $\theta$, so a feasible KKT solution is its unique minimizer.

The convergence criterion should remain the original projected log-rate
gradient, $\|P_\theta g_\theta\|_\infty<10^{-3}$. Applying the same numerical
threshold to a rescaled-coordinate gradient would silently change the test.

\section{Recommended optimization strategy}

\subsection{Short EM initialization, followed by observed-curvature learning}

The strongest currently supported combination is:
\begin{enumerate}
 \item Start from the same feasible rates as the baseline. Take two or three
       bounded EM steps using counts extracted in the existing gradient passes.
       Reuse the already-built model and parsed families.
 \item At the final warm-up endpoint $\theta_K$, form
       $C=C_\theta(\theta_K;N^{K-1})$ using the latest fixed count set and
       \eqref{eq:surrogate-curvature}. This is the exact Hessian of that
       fixed-count surrogate, not the observed Hessian at $\theta_K$.
 \item Calibrate its scale using the already-paid latest observed-gradient
       secant, $s=\theta_{K-1}-\theta_{K-2}$ and
       $y=g_{K-1}-g_{K-2}$. When both quadratic quantities are finite and
       positive, use
       \[
        \alpha=\frac{s^\tr y}{s^\tr Cs},\qquad B_0=\alpha C.
       \]
       Otherwise retain a safe positive scale. Apply a safeguarded BFGS
       correction on coordinates eligible at both secant endpoints.
 \item Continue with the original per-family trust-region/BFGS method,
       scheduled exact curvature refreshes, bound handling, and convergence
       tests. Keep the full $3\times3$ matrix: its algebra is cheap compared
       with a reconciliation gradient.
\end{enumerate}

For a positive-definite $B$ and suitable positive-curvature secant, the familiar
update is
\[
 B^+=B-\frac{Bss^\tr B}{s^\tr Bs}+\frac{yy^\tr}{s^\tr y}.
\]
Practical curvature floors and free-coordinate safeguards remain necessary.
The calibration learns only one direction and uses a secant preceding the
endpoint. It is a useful approximation, not a proof that $B_0$ equals observed
information or globally majorizes the objective.

Exact posterior counts and an exact bounded M-step imply EM likelihood
monotonicity. Pruned FP32 adjoints, finite fixed-point tolerances, and rounding
the endpoint back to model precision weaken that software guarantee. A double
output buffer does not turn FP32-derived counts into exact counts. Observed
likelihood checks and the unchanged final certificate remain important.

\subsection{Interpret adaptive curvature as local coordinates}

For a safeguarded positive-definite approximation $B_f$ in one family's free
subspace, write $B_f=R_f^\tr R_f$. Then
\begin{equation}
 z_f=R_f(\theta_f-\bar\theta_f)
 \label{eq:whitening}
\end{equation}
gives an identity quadratic curvature in the local model. These are
family-specific, adaptive coordinates: they can scale and rotate to follow the
posterior ambiguity. They need not be materialized in the likelihood kernels;
the step can be solved directly in $\theta$, with its simple box constraints.

A crucial distinction is that an SPD preconditioning metric transforms as a
tensor, while an ordinary off-stationary Hessian has the extra term in
\eqref{eq:hessian-transform}. Pulling the diagonal matrix
\eqref{eq:hierarchical-curvature} back without that term defines a possible
metric, but is not generally the exact fixed-count Hessian in log rates.
Computing \eqref{eq:surrogate-curvature} directly avoids this ambiguity.

Whitening alone is not an additional acceleration over solving the same
quadratic model exactly in the original variables. It helps when it improves
the approximation, numerical conditioning, or trust-region metric. Any changed
trust-region geometry should therefore be tested as its own algorithmic choice.

\subsection{Priorities for further experiments, not established gains}

\begin{enumerate}
 \item Prefer count-informed full-matrix curvature over a permanently diagonal
       hierarchical optimizer. The latter discards precisely the posterior
       coupling responsible for difficult families.
 \item Test any new scaling or metric with matched bounds, acceptance rules,
       and curvature-refresh schedules. A negative result for one hierarchical
       replay is not a theorem that every hierarchical hybrid fails.
 \item Consider improving the latent augmentation only if missing information
       remains the bottleneck. Collapsing avoidable latent variables or using a
       valid alternative augmentation can change EM's convergence spectrum;
       renaming the same parameters cannot. Such changes require a new
       likelihood/EM derivation, not ad hoc subtraction of ghost counts.
 \item Do not assume that a better median condition number or a shorter path
       to saved optima lowers runtime. Charge additional gradient evaluations,
       Hessian probes, rejected trials, initialization, and difficult-family
       completion. Saved fitted optima are evaluation references, never inputs
       to a proposed initialization.
\end{enumerate}

\section{Experimental evidence and its limits}

\subsection{Hierarchical coordinates improve geometry, but have not won the runtime comparison}

On 500 saved Coleman-family states, the hierarchical formulas, Jacobian, and
Hessian transformation were checked numerically against differentiation. At the
post-Adam point, observed-Hessian indefiniteness fell from $64\%$ in log rates
to $9\%$ in the hierarchy. The median relative curvature drift to the fitted
endpoint was still about $1.06$.

In the subsequent 200-family GPU replay, a scaled hierarchical diagonal reached
186/200 fresh certificates after the longer bounded run, with NLL about
$1.37$ bits worse than the fully converged production baseline. Some failures
were at coupled rate-bound corners. That replay lacked the production exact
refresh at step 15, so it was not an otherwise-identical full hybrid comparison.
The evidence supports rejecting that tested configuration, not ruling out every
coordinate-aware method.

\subsection{Count-informed initialization has measured gains}

Table~\ref{tab:local-results} gives the matched 500-family RTX 4090 prototype
runs recorded in the campaign. Both EM prototypes include their redundant
warm-model construction in wall time. All three use the same production
freeze-time projected-gradient criterion and certify all 500 families.

\begin{table}[htbp]
 \centering
 \small
 \begin{tabular}{lrrr}
  \toprule
  Initialization & Wall time (s) & Gradient clade equivalents & NLL (bits)\\
  \midrule
  Three Adam steps & 108.33 & 17.273 & 1618461.966967\\
  Two EM steps + scaled curvature & 91.51 & 13.660 & 1618461.994534\\
  Three EM steps + scaled curvature & 87.21 & 13.299 & 1618461.952767\\
  \bottomrule
 \end{tabular}
 \caption{Matched local prototype evidence; these are not full-Coleman H100 timings.}
 \label{tab:local-results}
\end{table}

The three-EM-step prototype is approximately $19.5\%$ faster than the Adam
prototype. Small aggregate NLL differences do not imply identical family-wise
solutions: four families change NLL by more than $0.01$ bits relative to Adam,
with both improvements and regressions. The recorded integrated three-step
gate subsequently took $81.42$ s on 500 families and certified all of them;
that number should not be presented as a matched integrated Adam speedup.

The first locally recorded full-Coleman H100 result is also encouraging:
two EM steps followed by the existing optimizer completed all 5,124 families
in $396.31$ s, with 28 Newton steps and all 5,124 freeze-time certificates.
This is about $23.9\%$ below the historical $520.5$ s record. Its fit-reported
NLL is $9049360.7436$ bits. However, the available contemporaneous Adam run
took $634.02$ s with unusually expensive startup/warm-up; it is not a fair sole
denominator for a warm speedup claim. The warmed paired control, repetitions,
and fresh common-model likelihood/gradient audits were not yet available in
the local result snapshot used here. The full-scale comparison is therefore
preliminary, not a completed paired validation.

For fair accounting, define gradient clade equivalents by
\[
 W_g=\sum_j
 \frac{\text{clades resident in the model evaluated on gradient call }j}
      {\text{clades in the initial full population}}.
\]
Frozen families still cost work until a replan removes them. Report Hessian
cost separately; its probes are not automatically equivalent to a fixed number
of gradient calls. The campaign's certificate uses the pruned FP32 gradient
recorded at freezing. It is not an unpruned mathematical stationarity proof.

\section{Conclusion}

There is a natural and exact hierarchical factorization of the event model.
It makes the fixed-count problem separable and provides useful geometric
diagnostics. The hard part of the observed problem, however, is the changing
posterior over histories, together with active bounds and expensive gradient
passes. A globally fixed coordinate transformation cannot erase that ambiguity.

\begin{quote}
Use complete-history structure to enter a good region cheaply and initialize
curvature; use observed-gradient information to learn each family's actual
local geometry; preserve the original constrained problem and its certificate.
\end{quote}

\section*{Provenance and scope}

This is an independent side-conversation note. It does not modify the earlier
mathematical report, the optimization code, or the running campaign. The evidence
is a snapshot of the local files read on 6 September 2026, not a continuously
updated campaign status. In particular, the older report's statement that EM
had not been measured is superseded by the later experiment records.

The model and earlier coordinate diagnosis are documented in
\path{agent_mathematical_report.tex}. Hierarchical identities, numerical checks,
and replay limitations are recorded in
\path{experiments/coleman_sol_20260906/geometry/PROGRESS.md}.
The survival-count and bounded-M-step review is in
\path{experiments/coleman_sol_20260906/geometry/EM_MATH_REVIEW.md}.
Local timing and likelihood comparisons are in
\path{experiments/coleman_sol_20260906/em/PROGRESS.md}.
The preliminary full-H100 result is
\path{experiments/coleman_sol_20260906/results/em_full_a.json}, with its
cold control in
\path{experiments/coleman_sol_20260906/results/baseline_interactive_a.json}.
Implementation details were checked against
\path{gpurec/fit/em_warmup.py} and the associated count API discussion.

\end{document}
